% Chapter 1

\chapter{Intro and Implementation Choices} % Main chapter title
\label{Chapter1} % For referencing the chapter elsewhere, use \ref{Chapter1} 
\setlength{\parindent}{0pt} % Rimuove il rientro di paragrafo

%----------------------------------------------------------------------------------------

% Define some commands to keep the formatting separated from the content 
\newcommand{\keyword}[1]{\textbf{#1}}
\newcommand{\tabhead}[1]{\textbf{#1}}
\newcommand{\code}[1]{\texttt{#1}}
\newcommand{\file}[1]{\texttt{\bfseries#1}}
\newcommand{\option}[1]{\texttt{\itshape#1}}

%----------------------------------------------------------------------------------------
\section{Smart Contract Implementation}

The decentralized lending service is implemented through a set of Solidity smart contracts developed and tested using the Hardhat framework. The main goal of the implementation is to provide a lending platform where contributors can supply funds, applicants can request loans backed by Bitcoin collateral information, and the community can participate in the approval process through a voting mechanism.

The smart contract architecture is composed of three main components: the \texttt{LendingService} contract, the \texttt{Loan} contract, and the \texttt{BitcoinOracle} contract. The separation of responsibilities between these contracts improves modularity and allows each component to manage a specific part of the protocol logic.

%----------------------------------------------------------------------------------------
\section{Lending Service Architecture}

The \texttt{LendingService} contract represents the core of the decentralized lending protocol. It manages the contributor liquidity pool, loan proposals, voting mechanism, compensation system, and the interaction with the oracle.

A main design choice was to separate the global lending logic from individual loans. Instead of storing all loan information inside a single contract, each approved proposal generates a new independent \texttt{Loan} contract. This approach allows every loan to maintain its own state, including the applicant, principal amount, repayment status, expiration block, and contributors involved.

The contract maintains the amount deposited by each contributor through the \texttt{deposited} mapping and tracks the funds currently allocated to active loans through the \texttt{locked} mapping. This distinction allows contributors to withdraw only the funds that are not committed to ongoing loans.

The deposit mechanism requires a minimum contribution defined by the \texttt{MIN\_DEPOSIT} constant. Contributors are stored inside an array to allow iteration during operations such as liquidity evaluation and vote weighting. When a contributor no longer has available funds or locked positions, the address is removed from the contributor list to reduce unnecessary storage usage.

%----------------------------------------------------------------------------------------
\section{Loan Proposal and Voting Mechanism}

Loan creation follows a proposal-based approach. Applicants submit a proposal containing the requested amount, interest rate, duration, and Bitcoin address used as collateral reference. Each proposal remains open for a fixed number of blocks, defined by the \texttt{PROPOSAL\_VOTING\_PERIOD} constant.

Only contributors with available deposited funds can participate in the voting process. Votes are weighted according to the contributor's disposable liquidity rather than using a simple one-address-one-vote mechanism. This choice prevents contributors with very small deposits from having the same influence as contributors providing a larger amount of liquidity.

When the voting period expires, the proposal is evaluated through multiple conditions:
\begin{itemize}
	\item the lending pool must contain enough disposable liquidity;
	\item the Bitcoin collateral value retrieved from the oracle must satisfy the liquidity requirement;
	\item the weighted approval votes must exceed the rejection weight.
\end{itemize}

If all conditions are satisfied, the service creates a new \texttt{Loan} contract and locks the required funds proportionally among the contributors.

%----------------------------------------------------------------------------------------
\section{Loan Contract Design}

The \texttt{Loan} contract represents a single lending agreement. It is deployed dynamically by the \texttt{LendingService} contract and receives the collected principal amount during construction. After deployment, the principal is immediately transferred to the applicant.

The contract stores immutable information about the loan, including:
\begin{itemize}
	\item applicant address;
	\item principal amount;
	\item interest rate;
	\item loan duration;
	\item collateral percentage;
	\item Bitcoin address associated with the collateral.
\end{itemize}

Using immutable variables for these parameters prevents modifications after deployment and reduces unnecessary storage operations.

The repayment mechanism allows partial or complete repayments. Each payment is divided proportionally into principal repayment and interest according to the configured interest rate. Principal repayments are distributed back to contributors, while interest payments are divided between contributor gains and the compensation pool.

The contributors participating in a loan are stored in a sorted order based on the initially locked amount. This ordering is used during repayment distribution, allowing refunds to be processed according to the initial contribution size.

%----------------------------------------------------------------------------------------
\section{Bitcoin Oracle Integration}

The lending protocol requires external Bitcoin information to evaluate whether a loan proposal has sufficient collateral. Since Ethereum smart contracts cannot directly access external blockchain data, a dedicated \texttt{BitcoinOracle} contract was implemented.

The oracle stores Bitcoin balances associated with Bitcoin addresses represented as byte arrays. Applicants can request an update by paying a minimum fee, while the oracle owner is responsible for pushing updated balances after performing the off-chain scan.

The separation between the request and update phases allows the blockchain component to remain independent from the off-chain Bitcoin data acquisition process. The \texttt{deploymentBlock} variable is also stored to support the off-chain service in identifying the starting point for synchronization.

%----------------------------------------------------------------------------------------
\section{Security and Administrative Choices}

The contracts include several mechanisms to improve security and maintainability.

The lending service implements an administrative mechanism based on a two-step ownership transfer. Instead of immediately changing the administrator, the current administrator nominates a pending administrator who must explicitly accept the role. This prevents accidental transfers to incorrect addresses.

A similar ownership transfer mechanism is implemented in the oracle contract.

The lending service also includes a termination mechanism that requires a successor contract address and prevents termination while active loans exist. This allows a future migration to a new contract version without leaving unresolved loans.

The implementation also includes callback functions between the \texttt{Loan} and \texttt{LendingService} contracts. These callbacks are used to notify the service about successful repayments, contributor refunds, interest collateral, and loan outcomes while keeping the two contracts logically separated.

Finally, the contract suite includes a vulnerable version of the lending service and a dedicated reentrancy attacker contract used during testing. These contracts are used to evaluate the impact of reentrancy attacks and validate the effectiveness of the adopted security modifications.

%----------------------------------------------------------------------------------------
